Deciding whether a multivariate polynomial of degree $2d \geq 4$ is nonnegative is known to be NP-hard.
However a sufficient condition for a polynomial to be nonnegative is easy to check.
We say that a polynomial is a \emph{sum of squares} (SOS) if there exist
polynomials $q_1, \ldots, q_M$ such that
\[ p(x) = \sum_{k=1}^Mq_k^2(x). \]
If a polynomial is SOS, then it is obviously nonnegative.

It is well known that if $p(x)$ is an homogeneous polynomial of degree
$2d$ then each $q_k(x)$ must be an homogeneous polynomial of degree $d$;
this can be shown easily using the Newton polytope of $p(x)$ and
\cite[Theorem~1]{reznick1978extremal}.  Let $x^{[d]}$ represent a
basis of the homogeneous polynomials of degree $d$.  We can check
whether a polynomial is SOS using semidefinite programming thanks to
the following theorem.
\begin{theorem}[\cite{choi1995sums,nesterov2000squared,parrilo2000structured,parrilo2003semidefinite,shor1987class}]
  \label{theo:sos}
  A homogeneous multivariate polynomial $p(x)$ of degree $2d$ is a sum
  of squares if and only if
  \[ p(x) = (x^{[d]})^\Tr Q x^{[d]} \]
  where $Q$ is a symmetric positive semidefinite matrix.
\end{theorem}

From the exact arithmetic viewpoint, the basis $x^{[d]}$ chosen in \theoref{sos} does
not affect whether $p(x)$ is SOS or not.  A specific choice of basis
may however improve the numerical behaviour of the corresponding
semidefinite program.

We denote the set of homogeneous polynomials of degree $2d$ as
$\Rh[x]$, the cone of homogeneous SOS polynomials of degree $2d$ as
$\Sosh$ and the dual of $\Sosh$ as $\Soshd$.

A common interpretation of the dual space $\Rhd$ of linear functionals on homogeneous polynomials of degree $2d$ is the space of moments of momonials of degree $2d$.
If $p(x) = a^\Tr x^{[d]}$ and $m$ is the vector of moments of $x^{[d]}$ of a measure $\mu$ then
\[ \la m, a \ra = \int p(x) \dif \mu = \la \mu, p \ra. \]
As a sum of squares polynomial is nonnegative,
this integral is nonnegative for any measure $\mu$.
Therefore, given a moment vector $m$, a necessary condition for a measure to exist with these moments is that $\la m, a \ra \geq 0$ for any vector of coefficients $a$ of a sum of squares polynomial.
That is, $\Soshd$ is a superset of the set of moments of measures.
The members of $\Soshd$ are often called \emph{pseudo-measures} and denoted $\pmu$; see~\cite{barak2012hypercontractivity}.
